\section{Discussion}
\label{sec:discussion}

Our results show that radio and X-ray observations are informative about recent cluster merger activity, while also demonstrating that the recoverability of this information depends on construction of the observable.
For individual simulated observations, the five-model ensemble achieves $R^2 = 0.203$ from the 3 Mpc radio images, and $R^2 = 0.406$ from the X-ray images.
These results extend image-based studies of cluster dynamical state toward continuous inference of a proxy for the merger history.
They also address observational limitations noted by \cite{Arendt2024}, which used idealized X-ray and Compton-y maps for classification.
Differences in targets, sample selection, and evaluation prevent a direct numerical comparison to their results.
\suggestion{Add the comparison requested in \texttt{todos.md}:
The ERGO-ML study of \citet{Chadayammuri2026} provides a complementary
view of the information in TNG-Cluster X-ray images. Their contrastive
representation of idealized maps captures variation in mass, dynamical
state, and merger history. Our supervised pseudo-TSC regression instead
tests recovery of a specified merger-history indicator, including from
individual forward-modelled observations. Differences in labels,
image construction, and evaluation prevent a direct accuracy comparison.
Their discussion of observational transfer motivates accounting for
effective area, detector gaps, and background noise, and testing
sensitivity to simulation physics. Our approximate forward model
addresses some of these effects but does not reproduce every observed
detector footprint or establish robustness across simulation suites.}

Properly interpreting the regression target is essential.
Pseudo-TSC is defined as the interpolated window size at which the merger score reaches a specific threshold.
Its strong rank correlation with the catalogue collision metric supports its use as an indicator of recent merger activity.
However, it is not calibrated as a direct time since pericentric passage.
The catalogue comparison includes minor mergers, with median mass ratio of approximately 1:11 for events defining the most recent collision time.
The 4 Gyr cap limits the proxy's age resolution, since the endpoint value can't distinguish histories with scores remaining below the threshold throughout the sampled window.
Predictions near the cap value are not intrinsically invalid, but cannot be interpreted as resolving collision times beyond that window.


The radio experiments indicate that both emission amplitude and spatial structure contribute to predictive accuracy.
Integrated emission is already a strong predictor of merger activity, and normalizing the radio maps to a common total emission reduces CNN performance.
Nevertheless, the normalized images retain substantial predictive information, as do the explicit morphological descriptors.
The CNN also remains predictive within mass bins, suggesting the presence of information in addition to the overall mass-age association.

However, our results do not provide a complete separation of information content between morphology, luminosity, and mass; they remain correlated within the simulated population, and the normalized-image predictions still correlate with the original luminosity.
Residual radio power at fixed mass decreases with pseudo-TSC in our simulated sample ($\rho = -0.33$), so brightness carries information about recent activity.
This is qualitatively consistent with time-dependent shock-powered radio light curves \citep{Nuza2024}.
However, its observational usefulness depends on whether this relation transfers to real clusters.

Before the density cut, a median of approximately 99\% of the radio weight arises in gas with $n_\text{H} > 10^{-4}$ cm$^{-3}$, including in a substantial cold component.
The cold, dense gas is consistent with stripped satellite material, but we did not investigate its origin.
Our shock emission model assumes non-radiative ICM shocks and neglects Coulomb losses \citep{HoeftBruggen2007}, motivating a check of the gas phases to which it is applied.
Our density cut is intended to restrict the model to its domain of applicability; our evaluation protocol demonstrates no clean-map accuracy gain after applying the cut ($R^2 = 0.500\to0.485$, paired-seed $p = 0.164$).
Sensitivity to the chosen gas selection remains to be quantified.

The radio amplitude is anchored to an observed radio-halo relation, although the simulated emission is powered by shocks, and the X-ray brightness correction is fitted using the observed sample.
Thus, agreement in overall brightness is partly imposed, rather than an independent validation of the emission model.
Interpreting residual differences requires comparisons between matched source definitions, mass distributions, apertures, and selection.

The adopted forward-modelled radio inputs yield substantially lower predictive scores than the clean-map benchmarks.
This comparison combines changes in resolution, aperture, background, and observing depth, and thus does not isolate the contribution from any one effect.
The improved performance of the 3 Mpc images is consistent with the presence of useful information outside the central cluster region.
These sensitivities reinforce the conclusion of \cite{Bottrell2019} that realistic image construction is important to determine whether features learned from simulations remain useful in observational conditions.

Radio and X-ray images provide complementary information in the clean-map experiments, but radio does not appear to provide additional benefit over X-ray alone after applying the forward-modelling pipeline.
However, the comparison uses a 1 Mpc radio aperture and different fold partitions, which somewhat limit its interpretation.
The apparent improvement from adding the 3 Mpc radio image to a noisy mass measurement is encouraging, but also requires a comparison with identical partitions.
\todo{should run 3Mpc box size, and also rerun with the fold partitions fixed}

The matched optical experiment provides clearer results, showing that, while idealized galaxy-structure features improve prediction beyond mass, redshift, and richness, adding them to X-ray images yields no demonstrated further gain.
These comparisons find no demonstrated incremental gain the tested observational band combinations.
However, since aperture, observing depth, and model choice limit the power of these comparisons, we can't conclude that the observationally-degraded bands contain completely identical information.

\cite{Lee2025} associate identified relics with particular merger events and construct a relation between relic separation and merger timing, including a strong relation for double-relics.
In contrast, our diagnostic measures the outermost thresholded peak in projected maps and compares its distance with the most recent catalogue collision.
Our outermost-peak diagnostic did not find a positive relation with the catalogue collision measure in the simulated sample.
This was a test of an inexpensive projected image statistic, rather than a reproduction of the three-dimensional relic identification and event association used by \citet{Lee2025}.
\todo{can we do this?}
A detected peak need not have been generated from the most recent merger event, especially when several collisions contribute to the radio morphology.
This result establishes that this simple peak statistic is ineffective for our sample.
Testing their relation would require matching the relic selection, merger association, and relevant sample cuts.

Several aspects of the observational model require approximations.
Image-plane radio injection, while including real background structure and imaging artifacts, does not reproduce the recovery of diffuse emissions using interferometric sampling and reconstruction.
Agreement in the morphological summary statistics of area fraction, perimeter, number of components, and Euler characteristic, does not establish full equivalence between mocks and real observations of all morphological indicators that may be used by the CNN, such as relative positions and orientations of components.

The X-ray forward model approximates the spectral response and background, but does not reproduce each observation's detector footprint, exposure structure, and roll angle.
Independent centering of the observed radio and X-ray images also removes their relative positional offsets.

The exposure experiment indicates that X-ray performance improves with greater depth, though with diminishing returns: the mean pooled-view $R^2$ increases from 0.312 at 5 ks exposure to 0.450 at 50 ks, then reaches 0.499 at 500 ks.

The systematic radio--X-ray offset in the LoVoCCS predictions (section~\ref{sec:lovoccs}) could arise from modality-dependent emission physics, brightness calibration, aperture and centering choices, or differences in regression bias.
A useful diagnostic would compare the two models' residuals for matched simulated observations at fixed mass, observing redshift, and pseudo-TSC, then test whether the observed offset is consistent with those biases.
\todo{would be pretty quick to run this probably}

While we were able to estimate some sources of uncertainty, others are more difficult to characterize.
The regression compresses the merger age range, systematically assigning older estimates to the youngest systems, and younger estimates to objects at the 4 Gyr age cap.
Scatter in predicted pseudo-TSC over different seeds measures variation between fitted networks.
Additional uncertainties arise from the finite halo sample, viewing direction, the mapping from pseudo-TSC to a particular collision, the emission model, and differences between simulated and observed selection and instrument response.
We have not shown these to be negligible; nevertheless, performance on the held-out sample demonstrates recovery of the merger history indicator, though uncertainties are not fully calibrated.
The distributional model's image-conditioned class probabilities outperform the baseline of fitting-sample frequency under categorical log score, ranked probability score, and no-crossing Brier score, despite providing no improvement in ensemble point prediction.
Calibration, coverage, and robustness under simulation-observation differences still require evaluation, especially for the no-crossing category.

Our exploratory negative results help identify where more work needs to be done to improve performance.
The frozen contrastive encoder experiment substantially underperforms the supervised control, but the cause has not been isolated.
Its positive pairs are rotations and reflections of the same observation, and the encoder is frozen during downstream training.
Alternative pair construction, objectives, or encoder fine-tuning would require controlled comparisons.
Thus, this pilot does not eliminate contrastive learning performing better under other conditions \citep{Chadayammuri2026}.

Using the diffusion model to generate synthetic data did not provide any performance improvement.
Because the data generation was not isolated from downstream test clusters, however, the comparison does not provide an independent estimate of the effect of augmentation, and the leakage does not establish why the observed difference was small.
\todo{can we make an independent estimate somehow?}
The comparison to CAMELS data combines a change in simulation domain with a change in merger-label definition, preventing attribution of its negative outcome to either factor.
These experiments motivate future work on more closely-matched experiments.

Next steps should therefore focus on making the inference more physically interpretable and observationally testable.
The radio prescription needs to be reconciled with the generated maps \todo{what does this mean?}, and sensitivity to the emitting-gas selection, fold assignment, and observing conditions should be established.
\suggestion{Resolve the prescription mismatch in the radio Methods
section before submission; it is not merely future work. The current
maps use a stored spectral-factor table and coefficient that differ
from the new analytical derivation. Either describe that implementation
or regenerate the affected results using the revised prescription.
Once resolved, replace this sentence with:
Sensitivity to the emitting-gas selection, fold assignment, and
observing conditions should be established before treating the inferred
proxy-age scale as robust across different modelling choices.}
Matching simulation outputs to observed selection and instrument response would strengthen conclusions about transferability, and consistent collision labels between simulation suites would test dependence on their differing physical models.
External dynamical constraints for real merging systems would allow the predicted ordering and timescale to be tested separately.
Weak-lensing and spectroscopic modelling of A3376 \citep{MonteiroOliveira2017}, for example, constrains a specific subcluster encounter independently of our image-based regression.
Such validation would require identifying whether the encounter corresponds to the history measured by our proxy, and matching the observing model to the validation data.
Additionally, independent simulation suits would provide a complementary test of sensitivity to simulation physics.
Such validation is necessary before image-based merger time estimates can be used quantitatively to study merger-driven particle acceleration.
